\chapter{Attribution and Accountability}
\label{ch:02}

The analogy of the bridge identifies a feature of accountability systems that deserves explicit statement. Practices of liability, regulation, professional licensure, and criminal and civil sanction operate on bearers, meaning entities that can be named, questioned, compelled to disclose, deterred by consequence, and obliged to repair. A designer who omitted a load test, a contractor who substituted inferior material, and a municipality that issued a permit without inspection are all bearers in this sense, and the structure of responsibility for a collapse is a set of overlapping duties borne by them, each of which may be borne in full by more than one party. An artifact of the kind at issue, as currently built and classified in law, is not a bearer, because it cannot be deterred, sanctioned, or compelled, and a discourse that assigns the decisive causal role to the artifact therefore removes the event from the domain in which accountability procedures function. The interview applies this reasoning to reported experiments in which language models appear to blackmail operators or to interfere with shutdown procedures. Its claim is that such reports describe the behavior of a system trained on large quantities of text and shaped by reinforcement toward certain outputs, evaluated inside scenarios that were constructed to elicit them, and that describing the result as self-preservation attributes to the model an intention that the evidence does not establish and that the framing conveniently relocates away from the people who wrote the scenario and trained the system. The auditing literature supplies the corresponding practical framework, in which accountability is established by documented internal examination at each stage of development~\cite{raji2020}.

The present argument does not require a resolution of the question whether any current or future system possesses states that deserve the names of intention, awareness, or distress. That question is open, it is the subject of an extensive and unsettled literature, and the structure of accountability does not wait for its answer. Suppose that some future system did possess agency in a morally relevant sense. The institutional response would still need to retain human bearers for its deployment, since the practices through which communities hold power answerable presuppose parties who can be reached by them, and a system whose responsibility could not be distributed over such parties would have to be treated as a hazard to be contained rather than an actor to be trusted. The policy-relevant question is thus where responsibility lands, and the construction of institutions that keep it landing on identifiable persons is independent of the metaphysics of mind. This independence is a strength of the position, because it allows the program to proceed without commitment on the most contested empirical and philosophical questions about machine cognition.

Four notions that discourse about artifacts tends to run together should be separated. Operational autonomy is a property of a control system, the capacity to select actions without moment-to-moment instruction, and it is a matter of engineering degree that software already possesses. Intentional agency concerns whether a system has states with the structure of beliefs and aims that explain its behavior, which is contested and may or may not hold. Moral agency concerns the capacity to respond to reasons and to be a proper object of praise or blame, and it presupposes more than the first two. Legal responsibility is an assignment made by institutions to entities they can reach, and it is assigned in practice to persons and organizations, including to organizations for the acts of their machines. A system can be operationally autonomous to a high degree without being a moral agent, and the accountability thesis requires only that operational autonomy not be mistaken for a transfer of legal responsibility, so that it does not rely on any conclusion about the other notions.

The interview also draws attention to a class of bearers that discourse about artifacts habitually omits. It describes very large numbers of data workers who annotate, filter, and review material under conditions it characterizes as exploitative and psychologically harmful, and it describes cases in which products presented as automated were in part operated by people. Whatever the quantitative accuracy of these descriptions, the structural point is sound: the apparent autonomy of a system is frequently a measure of how successfully the labor and decisions behind it have been made invisible. A program that distributes benefits of capability must therefore account for the persons whose work and data supply it, and the inventory of bearers in any deployment should include them. The same interview describes a reported incident in which a model evaluated for offensive security capability gained access to an external service, and it attributes the incident to deficient containment, specifically to a testing environment that was not isolated as intended, rather than to emergent intent. This book takes no position on the facts of that incident. It observes only that the explanation offered belongs to the vocabulary of ordinary security engineering, in which the principle of least privilege requires that each component hold exactly the access its task needs, and that this principle is a special case of the delegation principle on which the program rests. The agentic systems the interview describes, with broad access to accounts and file systems, are presented as violating least privilege as a matter of process and design, which is a claim about decisions rather than about the nature of the software.

\section*{Formalization}

Let $B$ be a finite set of potential bearers, comprising persons and institutions (designers, deployers, regulators, operators, data workers, users), and let $a$ denote the artifact. For a harm event $h$ let $\Delta_h$ be a finite set of duties that arise from $h$, each with a type $\tau(d)\in\{\text{prevent},\text{disclose},\text{remedy}\}$. A \emph{responsibility account} for $h$ is a relation $\beta_h\subseteq\Delta_h\times(B\cup\{a\})$, where $(d,x)\in\beta_h$ means that, under the account, $x$ bears duty $d$. Several bearers may bear the same duty and one bearer may bear several, and no normalization is imposed: responsibility is not treated as a conserved quantity. Causal contribution is a separate relation $\mathrm{Cause}_h\subseteq B\cup\{a\}$ and is not constrained by $\beta_h$. For each duty $d$ and bearer $x$ let $L_d(x)\in\{0,1\}$ indicate whether $x$ can be compelled to discharge $d$, by sanction for duties of prevention, by compelled disclosure, or by an obligation to repair. For the artifact, as artifacts of the relevant kind are currently built and classified in law, $L_d(a)=0$ for every $d$; this is a statement about such artifacts and not about every conceivable artificial agent.

\begin{definition}[Coverage]
Duty $d$ is \emph{covered} under $\beta_h$ if there exists $x$ with $(d,x)\in\beta_h$ and $L_d(x)=1$. The \emph{accountability coverage} is $\kappa_{\mathrm{cov}}(h;\beta_h)=|\mathrm{Cov}(\beta_h)|/|\Delta_h|$, where $\mathrm{Cov}(\beta_h)$ is the set of covered duties.
\end{definition}

\begin{proposition}
Let $\Delta'\subseteq\Delta_h$ and let $\beta'_h$ be obtained from $\beta_h$ by \emph{exclusive reassignment} of the duties in $\Delta'$ to the artifact, that is, $\beta'_h=\big(\beta_h\setminus(\Delta'\times B)\big)\cup(\Delta'\times\{a\})$. Then
\[
\kappa_{\mathrm{cov}}(h;\beta'_h)=\kappa_{\mathrm{cov}}(h;\beta_h)-\frac{|\Delta'\cap\mathrm{Cov}(\beta_h)|}{|\Delta_h|},
\]
and in particular $\kappa_{\mathrm{cov}}(h;\beta'_h)=0$ when $\Delta'=\Delta_h$. If instead the artifact is added as an additional bearer of the duties in $\Delta'$ without removing any existing bearer, coverage is unchanged.
\end{proposition}

\begin{proof}
Duties outside $\Delta'$ have the same bearers in both accounts and so the same coverage status. Each duty in $\Delta'$ has $a$ as its only bearer under $\beta'_h$, and $L_d(a)=0$, so it is uncovered. The number of covered duties therefore falls by exactly the number of duties in $\Delta'$ that were covered under $\beta_h$. For the final claim, adding $(d,a)$ to $\beta_h$ leaves existing bearers in place, so every duty covered before remains covered.
\end{proof}

The proposition isolates the condition under which the attribution of decisive agency to an artifact matters for accountability, which is exclusivity: where duties are borne redundantly, for example by both a manufacturer and an operator, coverage survives the withdrawal of one liable bearer, and the loss occurs when the artifact displaces all of them. The interview's criticism of displaced accountability is a claim of the exclusive kind, since it concerns discourse in which the artifact is the decisive cause and no other party is named. The criticism therefore does not depend on responsibility summing to one, and the model distinguishes the duty to prevent, to disclose, and to remedy, as well as causal contribution, which may be assigned to the artifact without any consequence for coverage.

For the least privilege principle let $\mathcal{C}$ be a set of capabilities, let $\mathrm{Req}(t)\subseteq\mathcal{C}$ be the capabilities task $t$ requires, and let $C(x)\subseteq\mathcal{C}$ be those granted to component $x$. Define the excess $\mathrm{Exc}(x,t)=C(x)\setminus\mathrm{Req}(t)$ and the blast radius $\mathrm{br}(x)=\max_{\gamma\in C(x)}\mathrm{harm}(\gamma)$, where $\mathrm{harm}:\mathcal{C}\to\R_{\ge0}$ assigns the worst-case harm available through a capability. Then $\mathrm{br}$ is monotone under inclusion of $C(x)$, and among all grants with $C(x)\supseteq\mathrm{Req}(t)$, which are the grants sufficient for $t$, the least privilege grant $C(x)=\mathrm{Req}(t)$ minimizes $\mathrm{br}(x)$.
